\documentclass[9pt, letterpaper]{extarticle}
\usepackage[utf8]{inputenc}
\usepackage[spanish]{babel}
\usepackage{hyperref}
\usepackage{titlesec}
\usepackage{enumitem}
\usepackage[margin=1in]{geometry}
\titleformat{\section}{\titlerule\vspace{7pt}\large\bfseries}{}{0em}{}
\titlespacing{\section}{0pt}{10pt}{5pt}

\begin{document}
\setlength{\parindent}{0pt}
    \begin{center}
        {\Huge \textbf{Jhon Gerardo Vega Medina}} \\ \vspace{6pt}
        Email: \href{mailto:jg.vegamd@gmail.com}{jg.vegamd@gmail.com} $\vert$
        Cel: (+57) 3102527117 $\vert$
        Bogotá, Colombia \\ \vspace{3pt}
        \href{https://www.linkedin.com/in/jhonvega}{www.linkedin.com/in/jhonvega} $\vert$
        \href{https://www.github.com/jhonvegam}{www.github.com/jhonvegam}
    \end{center}

    \section*{Perfil Profesional}
        \textbf{Analytics Engineer} con Maestría en Ingeniería Industrial (M.Sc.), enfocado en el diseño de arquitecturas de datos y modelado analítico para soportar decisiones de negocio a escala. Experto en \textbf{SQL}, \textbf{Python} y \textbf{BigQuery}, construyendo \textbf{pipelines} robustos e integrando fuentes como \textbf{CRM}, \textbf{GA4} y plataformas de marketing para definir métricas confiables y accionables.
    \vspace{5pt}
        He incorporado \textbf{Inteligencia Artificial} y \textbf{Machine Learning} en Marketing Operations para automatizar procesos y mejorar la eficiencia operativa. Me especializo en convertir datos complejos en modelos analíticos escalables que generan impacto en negocio en entornos \textbf{cloud (GCP, AWS, Azure)}.

    \section*{Experiencia}

        \noindent \textbf{Clio}
        
        \begin{itemize}[leftmargin=15pt, nosep]
            \item \textbf{Data Analyst Specialist - Marketing Operations} $\vert$ abril 2026 - ACTUALIDAD
            \begin{itemize}[leftmargin=15pt, label={--}, nosep]
                \item Lidero la integración de Inteligencia Artificial y modelos de Machine Learning para mejorar la eficiencia y escalabilidad de los procesos de la compañía.
                \item Desarrollo modelos predictivos basados en datos de CRM y comportamiento web para estimar la probabilidad de conversión de prospectos.
                \item Diseño e implemento agentes de IA para automatizar tareas repetitivas de alto volumen, incluyendo web scraping, entity matching y limpieza de datos.
                \item Optimizo flujos de trabajo mediante automatización, permitiendo al equipo enfocarse en tareas de mayor valor analítico.
            \end{itemize}
        \end{itemize}
        
        \vspace{10pt}

        \begin{itemize}[leftmargin=15pt, nosep]
            \item \textbf{Data Analyst - Marketing Operations} $\vert$ julio 2025 - marzo 2026 (9 meses)
            \begin{itemize}[leftmargin=15pt, label={--}, nosep]
                \item Lideré el desarrollo de productos de datos end-to-end, alineando necesidades de negocio con la arquitectura de datos para soportar decisiones estratégicas
                \item Diseñé y optimicé pipelines en BigQuery integrando HubSpot, GA4 y Google Ads, mejorando la disponibilidad de datos a inmediato.
                \item Implementé modelos de datos y métricas en SQL, garantizando consistencia y gobierno de la información.
                \item Desarrollé dashboards automatizados en Looker Studio para el seguimiento de KPIs de marketing, mejorando la visibilidad y toma de decisiones del equipo.
            \end{itemize}
        \end{itemize}

        \vspace{10pt}

        \noindent \textbf{Universidad de Los Andes}

        \begin{itemize}[leftmargin=15pt, nosep]
            \item \textbf{Asistente Graduado de Docencia} $\vert$ agosto 2023 - diciembre 2024 (1 año 6 meses)
            \begin{itemize}[leftmargin=15pt, label={--}, nosep]
                \item Enseñé fundamentos de Machine Learning en el curso de Probabilidad \& Estadística II, utilizando Python y R.
                \item Expliqué modelos supervisados (regresión lineal y logística) y no supervisados (clustering), aplicados a problemas reales.
                \item Traduje conceptos técnicos complejos en explicaciones accesibles para estudiantes con distintos niveles de experiencia.
                \item Apoyé el desarrollo de habilidades analíticas en estudiantes, facilitando la comprensión de técnicas de ciencia de datos.
            \end{itemize}
        \end{itemize}
        
        \vspace{10pt}

        \noindent \textbf{SUpost}

        \begin{itemize}[leftmargin=15pt, nosep]
            \item \textbf{Data Engineer - Consultor} $\vert$ agosto 2021 - octubre 2023 (2 años 3 meses)
            \begin{itemize}[leftmargin=15pt, label={--}, nosep]
                \item Implementé procesos ETL que mejoraron la disponibilidad y calidad de datos.
                \item Desarrollé soluciones de visualización que contribuyeron a la captación y retención de clientes.
            \end{itemize}
        \end{itemize}

    \section*{Educación}
        \textbf{Magister en Ingeniería Industrial (M.Sc.) - Enfoque en Data Analytics} \\
        \textit{Universidad de Los Andes} $\vert$ 2023 - 2024, Bogotá, Colombia
        \begin{itemize}[leftmargin=15pt, nosep]
            \item \textit{\textbf{Cursos relevantes:} Mastering machine learning, Analítica computacional para la toma de decisiones, Statistical learning for data analysis, Diseño de experimentos en ingeniería, Optimización avanzada.}
        \end{itemize}

        \vspace{10pt}

        \textbf{Ingeniero de Sistemas Computacionales} \\
        \textit{Universidad de Los Andes} $\vert$ 2017 - 2026 Bogotá, Colombia (en curso)

        \vspace{10pt}

        \textbf{Ingeniero Industrial} \\
        \textit{Universidad de Los Andes} $\vert$ 2014 - 2022 Bogotá, Colombia

    \section*{Habilidades}
        \textbf{Lenguajes \& Data Science:} Python (Pandas, Numpy, Scikit-learn, PyTorch), R, SQL.
        \textbf{Cloud \& Infraestructura:} GCP (BigQuery), Azure, AWS.
        \textbf{Marketing \& Analytics:} HubSpot, Salesforce, GA4, Google Ads, Looker Studio, Power BI.
        \textbf{IA \& Automatización:} Agentes de IA, Web Scraping.
        \textbf{Herramientas:} Git/GitHub. 
        
    \section*{Publicaciones}
        \textbf{Causal survival analysis for estimating the risk of renal failure in diabetic patients based on the health entity provider of the contributory regime to which they are affiliated in Colombia.} \\
        Repositorio de la Universidad de los Andes. (2024) \href{https://hdl.handle.net/1992/75234}{(Link Publicación)}
        
    \section*{Logros}
        Beca de asistencia graduada para cursar la Maestría en Ingeniería Industrial con enfoque en analítica de datos (2023-2024). Universidad de los Andes. (Bogotá, Colombia)

    \section*{Certificaciones}
        \textbf{Microsoft Certified:} Azure Fundamentals \href{https://learn.microsoft.com/es-es/users/jhonvega-6603/credentials/7704dbc5a72d05f?ref=https%3A%2F%2Fwww.google.com%2F}{(AZ-900)}, 2025

    \section*{Idiomas}
        \textbf{Ingles:} Profesional (Nivel B2) \\
        \textbf{Español:} Nativo

    \section*{Hobbies}
        Bailar, ajedrez, tenis de mesa, gimnasio

\end{document}